\chapter{Integrating Continuum Local Problems at the Material Boundary}
\label{ch:continuum-local-integrator-boundary}

\section{Motivation}

Chapter~\ref{ch:continuum-local-problem-interface} introduced a generic local
nonlinear problem driven by \texttt{ConstitutiveKinematics<dim>}. That was the
necessary abstraction step, but it was not yet sufficient for real use in the
library. A constitutive framework is only operational when the local problem
can be consumed through the same public boundary used by elements and material
sites:
\begin{itemize}
  \item \texttt{Material<Policy>},
  \item \texttt{MaterialPoint<Policy>},
  \item and, later, section-level constitutive sites built on the same handle.
\end{itemize}

The present phase performs that actual integration. The new
\texttt{ContinuumLocalProblemIntegrator<LocalProblem, LocalSolver>} is a
constitutive-integrator policy injected at the \texttt{Material<>} boundary.
This keeps the public constitutive API stable while opening a finite-strain
path that is still fully customizable at compile time.

\section{Architectural Position}

After this step, the continuum constitutive stack is:
\begin{enumerate}
  \item a constitutive relation
        \[
          \mathcal{R},
        \]
  \item an explicit algorithmic state
        \[
          \alpha_n,
        \]
  \item an injected local problem
        \[
          \mathbf{R}(\mathbf{x}; \mathcal{K}_{n+1}, \alpha_n) = \mathbf{0},
        \]
  \item an injected local nonlinear solver,
  \item an injected constitutive integrator at the erased
        \texttt{Material<>} boundary.
\end{enumerate}

In code, that last point matters. \texttt{Material<>} remains the uniform
polymorphic handle seen by the rest of the library, but it is no longer tied
to:
\begin{itemize}
  \item direct constitutive evaluation only,
  \item embedded return-mapping only,
  \item or small-strain kinematic carriers only.
\end{itemize}

The continuum-local path is therefore not a parallel API. It is a first-class
integration path through the existing constitutive boundary.

\section{Prediction and Commit Semantics}

This integration step also fixes an important semantic point.

For a rate-independent local constitutive update, evaluating the local problem
at a given kinematic state should generally be side-effect free:
\[
  \sigma_{n+1}^{\text{alg}}
  =
  \mathcal{G}\bigl(\mathcal{K}_{n+1}, \alpha_n\bigr),
\]
where $\mathcal{G}$ denotes the local constitutive solve.

That means:
\begin{itemize}
  \item \texttt{compute\_response(kin)} predicts the algorithmic response,
  \item \texttt{tangent(kin)} predicts the associated algorithmic tangent,
  \item \texttt{commit(kin)} persists the updated internal state and updates
        the site kinematic state.
\end{itemize}

This is a better contract than using \texttt{compute\_response()} as a hidden
mutation step. The distinction is especially important for:
\begin{itemize}
  \item repeated tangent evaluations inside a global Newton iteration,
  \item material-curve drivers,
  \item rollback/restart logic,
  \item future PETSc-based local solvers,
  \item and transfer to scale-reconstruction workflows.
\end{itemize}

In other words, the local constitutive solve is \emph{algorithmic}, but its
state persistence is still explicit.

\section{Mathematical Role of the Boundary Integrator}

Given:
\begin{itemize}
  \item a continuum carrier $\mathcal{K}_{n+1}$,
  \item a committed algorithmic state $\alpha_n$,
  \item a local problem with unknown $\mathbf{x}$,
\end{itemize}
the boundary integrator performs:
\begin{enumerate}
  \item build the local context
        \[
          \mathcal{C}_{n+1}
          =
          \mathcal{C}(\mathcal{K}_{n+1}, \alpha_n),
        \]
  \item solve the local nonlinear system
        \[
          \mathbf{R}(\mathbf{x}; \mathcal{C}_{n+1}) = \mathbf{0},
        \]
  \item reconstruct the constitutive result
        \[
          \bigl(
            \sigma_{n+1},
            \mathbb{C}_{n+1}^{\text{alg}},
            \alpha_{n+1}
          \bigr)
          =
          \mathcal{F}(\mathbf{x}, \mathcal{C}_{n+1}),
        \]
  \item on commit, store $\alpha_{n+1}$ and the corresponding kinematic state.
\end{enumerate}

The essential architectural point is that none of those steps are hard-coded
inside \texttt{Material<>}. The erased handle delegates them to the injected
integrator, which itself delegates to the injected local problem and local
solver.

\section{Compile-Time Extension Points Preserved}

The new path remains open at compile time in every important direction:
\begin{itemize}
  \item \textbf{Local problem}: different finite-strain models may define
        their own context, unknown vector, residual and reconstruction logic.
  \item \textbf{Local solver}: Newton is only one
        \texttt{LocalNonlinearSolverPolicy}; it can be replaced by damped,
        quasi-Newton, line-search or PETSc-backed solvers.
  \item \textbf{Constitutive integrator}: \texttt{Material<>} only depends on
        the integrator contract, not on the particular local-problem type.
  \item \textbf{Constitutive site}: the same integration path is available
        through \texttt{MaterialPoint<>}, which is the actual consumer at
        continuum Gauss points.
\end{itemize}

This keeps the direction aligned with the long-term goals:
\begin{itemize}
  \item finite-strain plasticity,
  \item damage-plasticity,
  \item fracture-oriented local updates,
  \item and multiscale transfer from structural to continuum models.
\end{itemize}

\section{Verification}

The regression in
\texttt{tests/test\_continuum\_local\_problem.cpp}
now checks the complete path:
\begin{itemize}
  \item direct use of
        \texttt{ContinuumLocalIntegrationAlgorithm<LocalProblem, Solver>},
  \item integration through
        \texttt{Material<ThreeDimensionalMaterial>},
  \item integration through
        \texttt{MaterialPoint<ThreeDimensionalMaterial>}.
\end{itemize}

The material-handle test is intentionally semantic:
\begin{itemize}
  \item prediction is verified to be side-effect free,
  \item commit is verified to persist the site kinematics,
  \item the committed response is checked against a reference constitutive site
        with the same explicit algorithmic state.
\end{itemize}

This is the right contract for the next step. The infrastructure is no longer
limited to an isolated continuum-local helper; it now participates in the real
constitutive boundary used by the rest of the codebase.

\section{References}

\begin{thebibliography}{9}

\bibitem{simo1998materialboundary}
J.~C. Simo and T.~J.~R. Hughes,
\emph{Computational Inelasticity},
Springer, 1998.

\bibitem{bonet2008materialboundary}
J.~Bonet and R.~D. Wood,
\emph{Nonlinear Continuum Mechanics for Finite Element Analysis},
2nd ed., Cambridge University Press, 2008.

\bibitem{desouza2008materialboundary}
E.~A. de Souza Neto, D.~Peri\'{c}, and D.~R.~J. Owen,
\emph{Computational Methods for Plasticity: Theory and Applications},
Wiley, 2008.

\end{thebibliography}
